\documentclass[10pt,a4paper]{amsart}
\usepackage{amsmath,amssymb,amsthm,graphicx}
\usepackage{verbatim}
\newenvironment{reference}{\comment}{\endcomment}
%\newenvironment{reference}{}{}
\newenvironment{slogan}{\comment}{\endcomment}
\newenvironment{history}{\comment}{\endcomment}

% For commutative diagrams we use Xy-pic
\usepackage[all]{xy}

% We use 2cell for 2-commutative diagrams.
\xyoption{2cell}
% Theorem environments.
%
\theoremstyle{plain}
\newtheorem{theorem}[subsection]{Theorem}
\newtheorem{proposition}[subsection]{Proposition}
\newtheorem{lemma}[subsection]{Lemma}

\theoremstyle{definition}
\newtheorem{definition}[subsection]{Definition}
\newtheorem{example}[subsection]{Example}
\newtheorem{exercise}[subsection]{Exercise}
\newtheorem{situation}[subsection]{Situation}

\theoremstyle{remark}
\newtheorem{remark}[subsection]{Remark}
\newtheorem{remarks}[subsection]{Remarks}

\numberwithin{equation}{subsection}

% Macros
%
\def\lim{\mathop{\mathrm{lim}}\nolimits}
\def\colim{\mathop{\mathrm{colim}}\nolimits}
\def\Spec{\mathop{\mathrm{Spec}}}
\def\Hom{\mathop{\mathrm{Hom}}\nolimits}
\def\Ext{\mathop{\mathrm{Ext}}\nolimits}
\def\SheafHom{\mathop{\mathcal{H}\!\mathit{om}}\nolimits}
\def\SheafExt{\mathop{\mathcal{E}\!\mathit{xt}}\nolimits}
\def\Sch{\mathit{Sch}}
\def\Mor{\mathop{\mathrm{Mor}}\nolimits}
\def\Ob{\mathop{\mathrm{Ob}}\nolimits}
\def\Sh{\mathop{\mathit{Sh}}\nolimits}
\def\NL{\mathop{N\!L}\nolimits}
\def\CH{\mathop{\mathrm{CH}}\nolimits}
\def\proetale{{pro\text{-}\acute{e}tale}}
\def\etale{{\acute{e}tale}}
\def\QCoh{\mathit{QCoh}}
\def\Ker{\mathop{\mathrm{Ker}}}
\def\Im{\mathop{\mathrm{Im}}}
\def\Coker{\mathop{\mathrm{Coker}}}
\def\Coim{\mathop{\mathrm{Coim}}}

% Boxtimes
%
\DeclareMathSymbol{\boxtimes}{\mathbin}{AMSa}{"02}

%
% Macros for moduli stacks/spaces
%
\def\QCohstack{\mathcal{QC}\!\mathit{oh}}
\def\Cohstack{\mathcal{C}\!\mathit{oh}}
\def\Spacesstack{\mathcal{S}\!\mathit{paces}}
\def\Quotfunctor{\mathrm{Quot}}
\def\Hilbfunctor{\mathrm{Hilb}}
\def\Curvesstack{\mathcal{C}\!\mathit{urves}}
\def\Polarizedstack{\mathcal{P}\!\mathit{olarized}}
\def\Complexesstack{\mathcal{C}\!\mathit{omplexes}}
% \Pic is the operator that assigns to X its picard group, usage \Pic(X)
% \Picardstack_{X/B} denotes the Picard stack of X over B
% \Picardfunctor_{X/B} denotes the Picard functor of X over B
\def\Pic{\mathop{\mathrm{Pic}}\nolimits}
\def\Picardstack{\mathcal{P}\!\mathit{ic}}
\def\Picardfunctor{\mathrm{Pic}}
\def\Deformationcategory{\mathcal{D}\!\mathit{ef}}
\usepackage{cleveref}
\crefname{originallemma}{original lemma}{original lemmas}
\Crefname{originallemma}{Original lemma}{Original lemmas}
\makeatletter
\newlabel{algebra-lemma-composition-flat}{{00HC}{source}{}{sourcepointer-1}{}}
\newlabel{algebra-lemma-composition-flat@cref}{{[originallemma][0][]00HC}{[1][0][]source}{}{}{}}
\makeatother
\crefname{originalremark}{original remark}{original remarks}
\Crefname{originalremark}{Original remark}{Original remarks}
\makeatletter
\newlabel{formal-spaces-remark-composition-variant-Noetherian}{{0GBK}{source}{}{sourcepointer-2}{}}
\newlabel{formal-spaces-remark-composition-variant-Noetherian@cref}{{[originalremark][0][]0GBK}{[1][0][]source}{}{}{}}
\makeatother
\crefname{originalreference}{original reference}{original references}
\Crefname{originalreference}{Original reference}{Original references}
\makeatletter
\newlabel{lemma-composition-finite-type}{{01T3}{source}{}{sourcepointer-3}{}}
\newlabel{lemma-composition-finite-type@cref}{{[originalreference][0][]01T3}{[1][0][]source}{}{}{}}
\makeatother
\crefname{originalreference}{original reference}{original references}
\Crefname{originalreference}{Original reference}{Original references}
\makeatletter
\newlabel{lemma-rig-flat-axioms}{{0GGU}{source}{}{sourcepointer-4}{}}
\newlabel{lemma-rig-flat-axioms@cref}{{[originalreference][0][]0GGU}{[1][0][]source}{}{}{}}
\makeatother
\title{Rig-flat maps of Noetherian adic rings are stable under composition}
\author{The Stacks Project Authors}
\date{stacks-project:restricted.tex:local-snapshot-20261012; GNU FDL 1.2 or later}
\begin{document}
\maketitle
\noindent\textit{Editorial note.} All delivered mathematical text is the author's own, copied from the pinned Stacks Project chapter file restricted.tex; no statement or proof step was rewritten, reordered or removed.\par\medskip
\noindent\textit{Editorial note.} The counted proof is the author's complete argument: the proof checks the composite is adic topologically finite type, localizes at each f, and composes the naive rig-flat maps using stability of flatness under composition\par\medskip
\noindent\textit{Editorial note.} The label lemma-composition-rig-flat-continuous names the counted statement, and the proof environment immediately following it in the source is delivered verbatim as its proof.\par\medskip
\noindent\textit{Editorial note.} The counted statement is carried by the author's own lemma environment in the chapter source.\par\medskip
\noindent\textit{Allowed standard prerequisites (editorial).}\begin{itemize}
\item the arrow property rig-flat on WAdm\textasciicircum{}Noeth (lemma-rig-flat-axioms of the same chapter) and the composition property of Formal Spaces, Remark formal-spaces-remark-composition-variant-Noetherian
\item composition of adic topologically finite type maps stays adic topologically finite type (lemma-composition-finite-type of the same chapter)
\item naive rig-flatness of the localizations A\_\{f\} -> B\_\{f\} and B\_\{f\} -> C\_\{f\} follows from rig-flatness
\item flat ring maps are stable under composition (Algebra, Lemma algebra-lemma-composition-flat)
\end{itemize}\medskip
\ifdefined\hypertarget\hypertarget{sourcepointer-1}{}\fi
\noindent\textit{Source reference.} Original lemma 00HC: sources/stacks-chow/tags.txt, source lines 638--638. Prerequisite pointer invoked at proof line(s) 4073; it names a standard fact of a sibling chapter that the counted argument uses, and every step of the counted proof built on it is retained verbatim in the delivered body.\par\medskip
\ifdefined\hypertarget\hypertarget{sourcepointer-2}{}\fi
\noindent\textit{Source reference.} Original remark 0GBK: sources/stacks-chow/tags.txt, source lines 20375--20375. Prerequisite pointer invoked at proof line(s) 4059; it names a standard fact of a sibling chapter that the counted argument uses, and every step of the counted proof built on it is retained verbatim in the delivered body.\par\medskip
\ifdefined\hypertarget\hypertarget{sourcepointer-3}{}\fi
\noindent\textit{Source reference.} Original reference 01T3: sources/stacks-restricted/pinned-original-source.tex, source lines 2533--2533. Prerequisite pointer invoked at proof line(s) 4067; it names a standard fact of the same chapter that the counted argument uses, and every step of the counted proof built on it is retained verbatim in the delivered body.\par\medskip
\ifdefined\hypertarget\hypertarget{sourcepointer-4}{}\fi
\noindent\textit{Source reference.} Original reference 0GGU: sources/stacks-restricted/pinned-original-source.tex, source lines 4010--4010. Prerequisite pointer invoked at proof line(s) 4063; it names a standard fact of the same chapter that the counted argument uses, and every step of the counted proof built on it is retained verbatim in the delivered body.\par\medskip
\begin{lemma}
\label{lemma-composition-rig-flat-continuous}
The property $P(\varphi)=$``$\varphi$ is rig-flat''
on arrows of $\textit{WAdm}^{Noeth}$ is stable under composition
as defined in Formal Spaces, Remark
\ref{formal-spaces-remark-composition-variant-Noetherian}.
\end{lemma}

\begin{proof}
The statement makes sense by Lemma \ref{lemma-rig-flat-axioms}.
To see that it is true assume we have rig-flat morphisms
$A \to B$ and $B \to C$ in $\textit{WAdm}^{Noeth}$.
Then $A \to C$ is adic and topologically of finite type
by Lemma \ref{lemma-composition-finite-type}.
To finish the proof we have to show that for all $f \in A$ the map
$A_{\{f\}} \to C_{\{f\}}$ is naively rig-flat.
Since $A_{\{f\}} \to B_{\{f\}}$ and $B_{\{f\}} \to C_{\{f\}}$
are naively rig-flat, it suffices to show that
compositions of naively rig-flat maps are naively rig-flat.
This is a consequence of Algebra, Lemma \ref{algebra-lemma-composition-flat}.
\end{proof}

\end{document}
